\clearpage
\SourceLeaf{445}
\fancyhead[L]{\small PERTURBATIONS OF PALLAS. POSTHUMOUS PAPERS.}
\pdfbookmark[1]{Posthumous papers. [II.] Exposition of a new method}{pallas-II-exposition}
\begin{center}
{\Large POSTHUMOUS PAPERS.\par}
\bigskip
{\large [II.] EXPOSITION\par}
{\large OF A NEW METHOD\par}
FOR CALCULATING PLANETARY PERTURBATIONS\par
WITH AN APPLICATION TO THE NUMERICAL CALCULATION\par
OF THE PERTURBATIONS OF THE MOTION OF PALLAS.
\end{center}
\begin{flushright}
\begin{minipage}{.94\textwidth}
\small\raggedright
. . . . . their motions, periods and their laws\\
Give me to scan.\quad \textsc{Thomsons} Autumn 1346.
\par\medskip
See, how associate round their central sun,\\
Their faithful rings the circling planets run\\
Exactly tracing their appointed sphere.\\
\hspace*{1em}Boyle on deity: Elegant Extracts p. 126.
\end{minipage}
\end{flushright}
\noindent\hspace*{\parindent}The calculation of planetary perturbations can take two different forms. In one form, the motion of a planet is referred to constant elements or to elements subject only to secular variations. The radius vector, longitude in the orbit, and latitude calculated on this assumption are then corrected by periodic equations. In the other form, the motion is referred to a variable ellipse, which represents the motion exactly at each instant without further correction. The success with which the first form alone has hitherto been used for the perturbations of the previously known planets is well known. It is recognized, however, that this success depends chiefly on several favorable circumstances that no longer hold in the same way for some of the new planets. The orbit of Pallas is inclined by 34 degrees to that of Jupiter and has an \SourceLeaf{446}eccentricity almost equal to \(\frac{1}{4}\), and the orbit of Juno, with a very considerable inclination, is still more eccentric. In these cases the old methods have appeared impracticable.
\par The new methods that form the subject of this memoir give the perturbations directly in the second form. It would indeed be easy to convert the perturbations of the elements, once found, into perturbations of the three coordinates if it were considered appropriate. In our view, however, this work can be dispensed with and the perturbations of the elements alone used. Although there are six of these elements and the perturbations of the three coordinates seem to require only half as much work in numerical applications, this disadvantage disappears whenever several positions of the planet separated by a modest interval of time are to be calculated. In that case, the most frequent in practice, a single set of elements will suffice, whereas the other method will require the inequalities of the coordinates to be calculated for at least two different epochs. The convenience of an auxiliary table for the equation of centre and the radius vector in a fixed ellipse must be given up when all the inequalities are referred to the elements themselves. This is also of very little consequence when working with a planet subject to as many inequalities as Pallas, and can even be offset by other simplifications in the calculation that it would be superfluous to recall here.
\par\bigskip\centerline{\rule{3cm}{.4pt}}\par
